\section{Method}
\label{sec:method}

\subsection{Frozen features under one interface}
\label{sec:features-interface}

Encoders differ in resolution, width and in how many frames they see. If each
gets its own probe, a difference in results could come from the probe. So every
encoder is reduced to the same input: its spatial feature grid, average-pooled to
$9\times16$ cells, projected on its first 64 principal components and scaled by
one global constant. The PCA is fitted on training logs only. The tokenizer
latents keep their native 16 channels. Table~\ref{tab:encoders} lists the
features.

\begin{table}[htbp]
\centering
\caption{Frozen features. Per-frame features are computed for each of the four
history frames, clip features once per scene from all four.}
\label{tab:encoders}
\begin{tblr}{colspec={lX[l]ll}, row{1}={font=\bfseries}, width=\textwidth}
\toprule
Name & Encoder & Frames & Channels \\
\midrule
$q_{\text{sem}}$, $q_{\text{rec}}$ & Orbis tokenizer, semantic and detail latents & per frame & 16 \\
DINOv2 & DINOv2 ViT-B/14~\cite{dinov2}, last layer & per frame & 64 of 768 \\
V-JEPA 2 & V-JEPA 2 ViT-L~\cite{vjepa2} on the 4-frame clip, newest time slot & clip & 64 of 1024 \\
WM$b$ & Orbis world model after block $b$, next-frame slot & clip & 64 of 768 \\
WM$b$-ctx & Orbis world model after block $b$, last observed frame & clip & 64 of 768 \\
\bottomrule
\end{tblr}
\end{table}

For the world model we run one forward pass over the four encoded history frames
and a pure-noise next frame ($\tau=1$), and keep the tokens after blocks 3, 7, 11,
15, 19 and 23. Two token positions are interesting. The next-frame slot is where
the model builds its prediction of the future. The last observed frame holds the
current scene as the model has contextualized it. We do not know in advance which
one carries more about the ego motion, so we test both.

\subsection{The probe}
\label{sec:probe}

\begin{figure}[htbp]
\centering
\begin{tikzpicture}[
  box/.style={draw, rounded corners=2pt, align=center, font=\small,
              minimum height=1.1cm, inner sep=4pt, text width=2.3cm},
  frozen/.style={box, fill=gray!15},
  arr/.style={-{Stealth[length=2mm]}, thick}]
\node[box, text width=1.6cm] (img) {4 front\\frames};
\node[frozen, right=0.5cm of img] (enc) {frozen encoder\\(Orbis, DINOv2,\\V-JEPA 2)};
\node[box, right=0.5cm of enc, text width=2.0cm] (grid) {$9\times16$ grid\\64 channels};
\node[box, right=0.5cm of grid, text width=2.9cm] (probe) {probe: 16 tokens per frame, query attends, 4 blocks};
\node[box, right=0.5cm of probe, text width=2.0cm] (out) {$K$ sampled\\trajectories};
\node[box, below=0.5cm of probe, text width=2.9cm, dashed] (ego) {ego status, command\\(optional tokens)};
\draw[arr] (img) -- (enc);
\draw[arr] (enc) -- (grid);
\draw[arr] (grid) -- (probe);
\draw[arr] (probe) -- (out);
\draw[arr] (ego) -- (probe);
\end{tikzpicture}
\caption{The probing setup. Only the probe is trained. It is a flow-matching
model on a 9-number polynomial code of the 4\,s trajectory.}
\label{fig:probe}
\end{figure}

The probe (Figure~\ref{fig:probe}) is a conditional flow-matching model in the
sense of Equation~\eqref{eq:fm-loss}. Each frame's feature grid is pooled by
attention into 16 tokens, and a learned embedding marks the frame. A single query
token, built from the noisy trajectory code and the flow time, attends to these
tokens through four transformer blocks with AdaLN-Zero time
conditioning~\cite{dit}, and predicts the velocity. For the ego-status
experiments the velocity and acceleration enter as one extra token and the
driving command as another.

\paragraph{Trajectory code.}
The flow does not act on the 24 raw numbers of the eight poses. Each of $x$, $y$
and $\psi$ is fitted by a cubic polynomial in time without a constant term (the
trajectory starts at the origin), which gives a 9-number code. The fit error on
training trajectories is 2.8\,cm. The code is smooth by construction and cannot
produce jagged paths. Section~\ref{sec:codec} compares it with raw poses.

\paragraph{Size.}
The probe has hidden width 16 and about 19k parameters. We kept it this small so
that it cannot learn much on its own and the results reflect the features.
Section~\ref{sec:capacity} shows what happens with a larger probe. All probes are
trained for 100k steps with batch 256 and a cosine learning-rate schedule from
$5\cdot10^{-4}$ to $5\cdot10^{-5}$, and sample with 10 Euler steps unless stated
otherwise.

\subsection{Metrics for a distribution of trajectories}
\label{sec:metrics}

For each scene we draw $K$ samples $X_1,\dots,X_K$ (16 during training, 32 in
offline evaluation) and compare them with the one logged future $y$. Let $d$ be
the average displacement between two trajectories (ADE).

\paragraph{Energy score.}
The main metric is the energy score~\cite{gneiting}, computed on the 16 numbers
$(x,y)$ of the eight poses:
\begin{equation}
\text{ES} = \mathbb{E}\,\lVert X - y\rVert - \tfrac12\,\mathbb{E}\,\lVert X - X'\rVert ,
\label{eq:energy}
\end{equation}
with $X, X'$ independent samples. It is a strictly proper scoring rule: its
expectation is minimized only when the samples follow the true conditional
distribution. The first term rewards samples close to the truth, the second
rewards spread. A collapsed sampler loses on the second term, a too-wide sampler
on the first. So, unlike ADE, the energy score cannot be gamed in either
direction.

\paragraph{Spread and dispersion.}
The spread is $\mathbb{E}\,d(X,X')$ and the dispersion is the ratio
\begin{equation}
\text{dispersion} = \frac{\mathbb{E}\,d(X,X')}{\mathbb{E}\,d(X,y)} .
\label{eq:dispersion}
\end{equation}
If the sampler is calibrated, $y$ is distributed like one more sample, so the
two expectations are equal and the dispersion is 1. Values below 1 mean the
samples are closer to each other than to the truth, which is the signature of an
overconfident or collapsed model. Values above 1 mean the samples are too wide.

\paragraph{Displacement errors.}
We also report the best-of-$K$ error $\text{minADE}_K=\min_k d(X_k,y)$, which
asks whether any sample covers the truth, and the error of the sample mean. The
NAVSIM literature reports the $L_2$ error at 4\,s of a single trajectory, which
for us is the sample mean.

\paragraph{EPDMS of a sampler.}
A planner must output one trajectory. We score either the first sample (one
random draw), the mean of the 32 samples, or the best of four samples per scene.
The last is an oracle and measures whether a good trajectory is among the
proposals.

\subsection{Reference points}
\label{sec:baselines}

Three baselines need no training. Constant velocity drives straight at the
current speed. The training mean predicts the average training trajectory for
every scene. The nearest-status lookup copies the future of the training scene
with the closest velocity and acceleration and the same driving command. On the
validation logs they reach ADE 2.76, 5.85 and 1.26\,m. The lookup halves the
error of constant velocity without seeing an image, which confirms that the ego
status is a strong shortcut on this data~\cite{egostatus}. We also train the probe
with no input at all (it learns the marginal distribution of trajectories) and
with the ego status and command only.

Before working on real data we checked the flow-matching denoiser on synthetic
2D curves whose distribution is known in closed form. It reproduced the fan of
target curves, and per-channel normalization of the targets was needed to avoid
wavy and overshooting samples. For the same reason every dimension of the
trajectory code is standardized with training statistics on real data.
